\chapter{BCS Project Criteria and Self-Reflection}
\label{ch:bcs}

% BUDGET: 3 pages. First person. At least one full page of reflection.

\section{The BCS project criteria}

\begin{description}[style=unboxed,leftmargin=0pt,itemsep=0.4em]

\item[An ability to apply practical and analytical skills.] The practical skills
are in Chapter~\ref{ch:design}: seven agents, each a graph with a validated
input/output contract, a sandboxed execution path for model-generated code, a
failure taxonomy that routes three classes of failure without a model call at all,
and deployment on a shared HPC facility under queue limits, pre-emption and a
storage quota. The analytical skills are in Chapter~\ref{ch:evaluation}, and the
work I would point to is not the building but the measuring: establishing that
single-run comparison on the benchmark is unsound because identical inputs produced
interpretable scores of 4, 6, 3, 2, 3 and 2, and then holding to that when a single
run would have told a better story.

\item[Innovation and creativity.] The pipeline's structure --- every agent a graph
rather than only the orchestrator --- is a design choice made so that each
deterministic check is an addressable, inspectable node rather than a line inside a
function. The preference dataset of Chapter~\ref{ch:po} is the contribution I
think is most genuinely novel: the fix loop already produces a rejected generation,
the check that rejected it and an accepted replacement, which is exactly a
preference triple, and the label comes from a compiler or an exit status rather
than from a human or a model's opinion.

\item[Synthesis of information, ideas and practices, with evaluation.] The
requirements analysis in Chapter~\ref{ch:requirements} turned two staff interviews
into six traceable design decisions, one of which (D1) overturned the persona-based
architecture my own proposal specified. The literature positioned the work against
AutoMETA and the autonomous-research-loop systems, and Chapter~\ref{ch:evaluation}
evaluates the result at four levels --- component tests, a frozen benchmark, an
ablation and three end-to-end sweeps --- including the results that do not flatter
it.

\item[Meeting a real need in a wider context.] The need came from the two
respondents directly: both said they would not accept unchecked output, and
Respondent~B's objection was specifically about opacity. That is why verification,
not automation, became the centre of the design. The wider intent is the one I
began with --- to let researchers who cannot program produce meaningful
computational work --- and Chapter~\ref{ch:ethics} sets out both sides of that,
including the risk that a tool of this kind lowers the cost of producing plausible
research that nobody has checked.

\item[Self-management of a significant piece of work.] The project ran across a
single term against a fixed deadline, on a cluster I do not control. I planned in
sweeps rather than features: build, run a batch, read what it did, repair, run the
same questions again. Three full sweeps of thirty-six questions were completed on
the same questions, model and cluster, which is what makes
Table~\ref{tab:three-sweeps} a comparison rather than an anecdote.

\item[Critical self-evaluation.] The section below, and --- more to the point --- the
fact that Chapter~\ref{ch:evaluation} leads with the project's failures, several of
which I introduced myself.

\end{description}

\section{Reflection}

The single most useful thing I learned has an uncomfortable shape: \emph{a
mechanism can be correct in the source file and do nothing in the run}. I found
this three times, and only the third time did I go looking for it.

The first was the citation checker, which reported 3{,}514 fabricated citations
across 108 drafts when there were four. It was scanning the paper's own reference
list. The second was a confirmation gate, written precisely to stop the system
trusting a dataset found by keyword search, which fired zero times across a
thirty-six-question sweep because the field it tests was set by one code path and
not by the path that resolves most data in production. The third is the one I find
hardest to write down. To make a benchmark case find a text corpus, I staged an
alias filename packing every phrasing the system might use --- ``dataset'',
``benchmark'', ``standard'', ``collection'' --- against a matcher that accepted a
single shared word. It worked, for that case. It also answered twenty-eight of
thirty-six requirements in the next sweep, and across every run in this project
nineteen experiments published a verdict computed on a file that could not answer
their question. A requirement naming MIMIC-III clinical records was answered with
newsgroup posts.

What makes that the most instructive failure is not the bug. It is that I made it
while trying to improve the system, that it passed every test, and that the
verdicts it produced looked entirely ordinary --- nine \emph{supported}, ten
\emph{refuted}, distributed across plausible questions. Nothing downstream could
have caught it, because a staged file is real data a human put there, which is the
one case the whole provenance design treats as trustworthy. A convenience added to
make one case pass was a change to the system, and I did not treat it as one.

The related lesson is about instruments. I lost real results to a verdict recorded
as the string \texttt{"True"} rather than a boolean, and twice built a
preference-accuracy metric that returned a clean, believable null --- once because
it averaged over the prompt as well as the completion, once because it was
measuring sentence length. In each case the measurement was wrong in a way that
looked like a finding. I now distrust a result that arrives tidy, and the habit I
would keep from this project is to predict what a metric should read when nothing
is happening, and check that first.

Testing deserves an honest verdict too. The requirement that all non-model logic be
unit-testable without a model present (NFR5) was worth the effort: 1{,}351 tests
run with no network, no model and no cluster, and the regression tests are written
from real failures rather than imagined ones. But the suite never caught any of the
three failures above, because each component did exactly what its tests asked. What
caught them was reading what a run actually produced. If I started again I would
budget as much time for instrumenting sweeps as for writing tests, and I would
measure a guardrail's coverage --- how often it fires, on what --- as a first-class
property alongside its logic.

On the proposal's reinforcement learning component, I still think reframing it was
right, and Chapter~\ref{ch:po} argues why: for a failure whose absence can be
checked mechanically, verification beats optimisation, because it is exact, works
on the first occurrence, and fails visibly. What I would concede is that I reached
that argument by building the verification first and justifying it afterwards,
rather than by reasoning it through at the proposal stage.

Two things I would do differently are concrete. I would stage real datasets at the
start rather than discovering, most of the way through, that the ceiling on
interpretable results was a data-availability problem rather than a code-generation
one. And I would treat the Experiment Planner as a source of defects much earlier:
by the end I had three gates protecting against bad data and bad compute, and then
found an experiment whose design was circular --- predicting a median split of
entropy from entropy --- which no amount of correct code could have rescued.

Finally, the thing I am most glad I did not do. Every one of the repairs in
Chapter~\ref{ch:evaluation} makes the system report \emph{less}: fewer verdicts,
more inconclusive results, and a last sweep that published nothing at all. It was
tempting, more than once and particularly with a deadline approaching, to loosen a
gate and let the numbers look better. The project is more useful for having
resisted that, and a system that says ``I cannot tell you'' is the only kind I would
be willing to put in front of a researcher who cannot check its work themselves.
